% Generated draft: requires mathematical and source review.
\documentclass[11pt]{article}
\usepackage[margin=25mm]{geometry}
\usepackage{fontspec}
\setmainfont{Latin Modern Roman}
\usepackage{amsmath,amssymb,mathtools,unicode-math}
\setmathfont{Latin Modern Math}
\usepackage{enumitem,needspace,xurl,hyperref}
\setlength{\parindent}{0pt}
\setlength{\parskip}{5pt}
\setlength{\emergencystretch}{3em}
\newcommand{\N}{\mathbb{N}}
\newcommand{\Z}{\mathbb{Z}}
\newcommand{\Q}{\mathbb{Q}}
\newcommand{\R}{\mathbb{R}}
\newcommand{\C}{\mathbb{C}}
\newcommand{\F}{\mathbb{F}}
\newcommand{\A}{\mathbb{A}}
\newcommand{\PP}{\mathbb{P}}
\newcommand{\E}{\mathbb{E}}
\newcommand{\eps}{\varepsilon}
\newcommand{\dd}{\,\mathrm{d}}
\begin{document}
\section*{UnsolvedMath ID: 30005675}
\textbf{Source problem code:} catalogue entry
\subsection{Definitions and mathematical statement}
Let \(G\) be a compact Lie group, and let \(\mathcal M_{\mathrm{FG}}^G\) denote the equivariant formal-group moduli stack intended by the supplied source record. The record does not supply its construction or a comodule presentation. For the intended presentation, define \(\mathcal C_G\) to be the full category of sheaves corresponding to comodules whose underlying modules are finitely presented. Here finite presentation means having finitely many generators and finitely many relations, with any grading conventions taken from the intended presentation. The coefficient ring, comodule structure, and grading conventions are not specified in the supplied record.

\textbf{Conjecture (supplied corrected formulation).} For every compact Lie group \(G\), the category \(\mathcal C_G\) is abelian. This is the meaning of coherence asserted for \(\mathcal M_{\mathrm{FG}}^G\) in the record. Explicitly, \(\mathcal C_G\) is to be additive, every morphism \(f\colon M\to N\) is to have a kernel and a cokernel in \(\mathcal C_G\), and the canonical morphism
\[
\operatorname{coker}(\ker f\longrightarrow M)
\longrightarrow
\ker(N\longrightarrow\operatorname{coker}f)
\]
is to be an isomorphism.

\textbf{Unresolved source qualification.} This formulation follows the supplied statement and clean statement. The supplied original excerpt instead states a conjecture about universal support and a Balmer-spectrum homeomorphism for finite \(G\)-spectra; it does not state the coherence conjecture. Thus the excerpt alone does not establish the intended stack, its comodule presentation, or the scope of the coherence assertion. In particular, the quantifier over all compact Lie groups is retained from the corrected formulation and requires source verification.
\subsection{Short English statement}
The supplied conjecture asks whether, for every compact Lie group \(G\), the comodules representing finitely presented sheaves on the intended equivariant formal-group stack form an abelian category. The stack and presentation are not defined in the record, and the original excerpt concerns a different conjecture. This is a statement-only draft requiring mathematical and source review.
\subsection{Sources}
\begin{itemize}
\item Supplied local UnsolvedMath registry record, permanent ID \(30005675\), source problem code \texttt{catalogue entry}, titled \emph{Coherence of Equivariant Formal-Group Moduli Stacks}; collection: Oberwolfach Reports Open Problems. No proposer is supplied; the proposed year is recorded as 2023.
\item Supplied report citation: \emph{Homotopy Theory (2024)}, Oberwolfach Reports, DOI \texttt{10.4171/owr/2023/34}, pp.\ 1893--1952. \url{https://doi.org/10.4171/owr/2023/34}. The record labels its corrected statement \texttt{corrected\_verified} and says it was reconstructed and checked against the official report. These are supplied verification claims, not independent verification here.
\item Supplied literature reference, labeled Hausmann--Meier: \url{https://arxiv.org/abs/2309.00850}. No title or more detailed attribution is supplied.
\item Literature assessment dated exactly \texttt{2026-08-22}: the supplied curation labels the coherence conjecture open and attributes the universal support theorem and Balmer-spectrum homeomorphism for compact abelian Lie groups to Hausmann and Meier. This is supplied dated curation, not an independently verified assessment. Its listed sources are \url{https://doi.org/10.4171/owr/2023/34} and \url{https://arxiv.org/abs/2309.00850}. The record explicitly requests independent verification of time-sensitive status and cited claims.
\end{itemize}
\end{document}
